\section{The complete low-degree head}\label{sec:head}
\subsection{Squared balance on the target coordinates}
\begin{lemma}[Active-coordinate squared balance]
\label{lem:positive-active}
If $P$ has total Walsh degree at most $r$, then
\begin{equation}
 \mathbb E(fP)_-^2
 \ge\left(\frac12-\sqrt{\frac rk}\right)\|P\|_2^2.
 \label{eq:positive-active-balance}
\end{equation}
For every fixed $\delta>0$, there is a universal $A_\delta$ such that
if $r\le k/64$ and $n\ge A_\delta D_r\log(D_r+1)$, then, with
probability at least $1-\delta$,
\begin{equation}
 \mathbb P_n(fP)_-^2\ge\frac14\|P\|_2^2
 \quad\text{for every polynomial $P$ of degree at most $r$}.
 \label{eq:positive-active-sampling}
\end{equation}
\end{lemma}
\begin{proof}
For $j\in U$ put $D_jP=(P-P\circ\mathrm{flip}_j)/2$.
Because $f$ changes sign under these flips, the inequality
$|F(s)-F(t)|\le(|s|+|t|)|s-t|$ for $F(t)=t|t|$ gives
\[
 |\mathbb E fP|P||\le2\|P\|_2\|D_jP\|_2.
\]
Averaging over $j\in U$ and using
$\sum_{j\in U}\|D_jP\|_2^2\le r\|P\|_2^2$ proves
\eqref{eq:positive-active-balance} from
$(fP)_-^2=(P^2-fP|P|)/2$.
The vector of all Walsh characters of degree at most $r$ is isotropic
and has squared length $D_r$ at every input. Under $r\le k/64$,
the population bound is at least $3\|P\|_2^2/8$.
Apply Lemma~\ref{lem:squared-sampling} at accuracy $1/8$ and use
homogeneity to obtain \eqref{eq:positive-active-sampling}.
\end{proof}


\subsection{Separating the first omitted pair of Fourier levels}
\begin{lemma}[The first paired shell]
\label{lem:positive-first-shell}
Fix $A\ge1$, $\delta>0$, and an odd $r_0\ge11$.
There is a sufficiently small $\alpha>0$ such that, uniformly for all
large $d$, the following holds. Suppose $r$ is odd and
\[
 r_0\le r\le\alpha d,\qquad
 A D_r\log(D_r+1)\le n<A D_{r+2}\log(D_{r+2}+1).
\]
Write $\rho=r/D$, $H=\log(e/\rho)$, and $\tau=\tau_r$.
Then, with probability at least $1-2\delta$,
\begin{equation}
 \|(Q_r)_{SS}\|_{\mathrm{op}}\le C_{L,A,\delta}H\tau.
 \label{eq:positive-improved-gram}
\end{equation}
Deterministically,
\begin{align}
 n\lambda_{r+1}&\le C_L A H\tau,
 \label{eq:positive-first-eigen}\\
 n\lambda_{r+3}&\le C_L A\rho^2H\tau.
 \label{eq:positive-second-eigen}
\end{align}
\end{lemma}
\begin{proof}
Let $B_r=K_{(r,r+2]}$, so $Q_r=B_r+Q_{r+2}$.
The two levels in $B_r$ have the common normalized eigenvalue
$\ell=\lambda_{r+1}=\lambda_{r+2}$ and total dimension
\[
 m=\binom d{r+1}+\binom d{r+2}=\binom D{r+2}.
\]
Its diagonal is $\Delta=m\ell=\tau_r-\tau_{r+2}$.
Choose the feature vector
$v(x)=\sqrt\ell(\chi_V(x))_{r<|V|\le r+2}\in\mathbb R^m$.
Then $\|v(x)\|_2^2=\Delta$ and
$\mathbb E vv^{\mathsf T}=\ell I_m$.
For $A_i=v(X_i)v(X_i)^{\mathsf T}-\ell I_m$,
\[
 \|A_i\|_{\mathrm{op}}\le\Delta,\qquad
 \left\|\sum_{i=1}^n\mathbb E A_i^2\right\|_{\mathrm{op}}
 \le n\ell\Delta.
\]
The self-adjoint matrix Bernstein inequality
\cite[Theorem~1.4]{tropp2012}, with $h=\log(2m/\delta)$, therefore gives,
with probability at least $1-\delta$,
\begin{equation}
 \|(B_r)_{SS}\|_{\mathrm{op}}
 \le n\ell+C\sqrt{n\ell\Delta h}+C\Delta h.
 \label{eq:positive-band-bernstein}
\end{equation}
Here the nonzero eigenvalues of the sample Gram matrix equal those of
$\sum_i v(X_i)v(X_i)^{\mathsf T}$.

The sharp spectrum and maximality interval imply
\[
 n\le C A m rH,\quad
 \ell\le\frac{C_L\tau}{rm},\quad
 \Delta\le\frac{C_L\tau}{r},\quad
 h\le C_\delta rH.
\]
Thus every term in \eqref{eq:positive-band-bernstein} is at most
$C_{L,A,\delta}\tau H$.
The fine-tail conclusion of Lemma~\ref{lem:sparse-gram} for the current
adaptive cutoff $r$ gives $\|(Q_{r+2})_{SS}\|\le2\tau$ with probability
at least $1-\delta$
after reducing $\alpha$ and increasing $d_0$.
The triangle inequality proves \eqref{eq:positive-improved-gram}.
The first eigenvalue estimate is among the displayed spectral bounds;
the two-step binomial ratio
$\binom D{r+2}/\binom D{r+4}\le C\rho^2$ proves the second.
No independence between the two sample events is used.
\end{proof}


\subsection{The complete head at an exact sample size}
\begin{proposition}[Head norm and population leakage]\label{prop:head}
Let $r$ be odd, $r\le k/64$, $\rho=r/(d+1)$,
$H=\log(e/\rho)$, and $n=n_r$ from~\eqref{eq:exact-size}.
Outside an event of arbitrarily small fixed probability, every minimizer
of~\eqref{eq:obstacle}, with $P=z_{\le r}$, satisfies
\begin{align}
 \|z\|_K^2&\le C n a^2d^\beta,\label{eq:budget}\\
 \|P\|_2^2&\le C a^2\rho^{-\beta}H,\label{eq:head-norm}\\
 \|z-P\|_2^2&\le C a^2\rho^{2-1/400-\beta}H.
 \label{eq:head-leakage}
\end{align}
The estimates hold simultaneously for all minimizers.
\end{proposition}
\begin{proof}
The interpolation comparator gives~\eqref{eq:budget}; at the displayed
sample size, its $d$ term is absorbed by $nd^\beta$.
The one-sided squared sampling theorem and the parity balance estimate
give
\[
 \frac14\|P\|_2^2\le\frac1n\sum_i(f(X_i)P(X_i))_-^2
 \le\frac1n\sum_i(z(X_i)-P(X_i))^2
 \le\frac{\|z\|_K^2}{n}\|(Q_r)_{SS}\|.
\]
Lemma~\ref{lem:positive-first-shell} bounds the last Gram matrix by
$CH\tau_r$, proving~\eqref{eq:head-norm}.
The first omitted pair has common eigenvalue
$\lambda_{r+1}=\lambda_{r+2}$, dimension
$\binom{d+1}{r+2}\asymp D_r\rho^{-2}$, and trace at most
$C\tau_r/r$. Therefore
\[
 n_r\lambda_{r+1}\le C\tau_r H\rho^{2-1/400}.
\]
Since $\|z-P\|_2^2\le\lambda_{r+1}\|z\|_K^2$, this proves
\eqref{eq:head-leakage}. The sample size belongs to the interval in
Lemma~\ref{lem:positive-first-shell}, because $1/400<2$.
\end{proof}

In particular the low-degree head inherits the sign of $z$ outside a
set of vanishing population measure, and its mass in $(-a/2,a/2)$ tends
to zero. The estimate~\eqref{eq:head-norm} concerns the entire head;
the much weaker bound $\lambda_0\|z\|_K^2$ would not suffice below.
